% 核心结论：两条高斯定律使用散度定理，两条环路定律使用斯托克斯定理。
% 图契约：物理竞赛讲义；印刷宽16cm，最小字号8.5pt；几何示意与公式分区。
% 布局：[几何/曲线 7cm] [定义/结果 7cm]；实线主量，虚线辅助；无装饰数据。
\documentclass[tikz,border=6pt]{standalone}
\usepackage{fontspec,xeCJK,amsmath,amssymb,graphicx}
\IfFontExistsTF{Times New Roman}{\setmainfont{Times New Roman}}{\setmainfont{TeX Gyre Termes}}
\IfFontExistsTF{Heiti SC}{\setCJKmainfont[AutoFakeBold=1.5]{Heiti SC}}{\setCJKmainfont[AutoFakeBold=1.5]{FandolHei-Regular}}
\usetikzlibrary{arrows.meta,bending,calc,positioning,patterns}
\definecolor{ink}{HTML}{243745}
\definecolor{blue}{HTML}{275DAB}
\definecolor{green}{HTML}{177C63}
\definecolor{red}{HTML}{BD4935}
\definecolor{orange}{HTML}{AA6C16}
\definecolor{purple}{HTML}{715393}
\tikzset{
 every node/.style={font=\fontsize{9.5}{13}\selectfont,text=ink},
 axis/.style={-{Stealth[length=4pt]},line width=.65pt,draw=ink!65},
 vec/.style={-{Stealth[length=5pt]},line width=1pt,draw=blue},
 short/.style={-{Stealth[length=3pt]},line width=.8pt},
 guide/.style={dashed,draw=ink!40,line width=.5pt},
 card/.style={rounded corners=2pt,draw=blue!60,fill=blue!4,
   align=center,text width=6.8cm,inner xsep=9pt,inner ysep=7pt},
 title/.style={font=\fontsize{11}{14}\selectfont\bfseries,text width=15cm,align=center},
 label/.style={font=\fontsize{8.5}{11}\selectfont,align=center}}
\begin{document}
\begin{tikzpicture}[x=1cm,y=1cm]
\path[use as bounding box] (-8,-3.35) rectangle (8,3.35);
\node[title] at (0,2.95) {麦克斯韦方程：积分形式与微分形式};

\node[card,text width=6.55cm] at (-4,1.2) {电场高斯定律\\[4pt]$\displaystyle\oint_{\partial V}\mathbf E\cdot d\mathbf S=\frac1{\varepsilon_0}\int_V\rho\,dV$\\[4pt]$\nabla\cdot\mathbf E=\rho/\varepsilon_0$};
\node[card,text width=6.55cm,draw=green!65,fill=green!4] at (4,1.2) {磁场高斯定律\\[4pt]$\displaystyle\oint_{\partial V}\mathbf B\cdot d\mathbf S=0$\\[4pt]$\nabla\cdot\mathbf B=0$};
\node[card,text width=6.55cm,draw=orange!65,fill=orange!4] at (-4,-1.25) {法拉第定律（固定曲面）\\[4pt]$\displaystyle\oint_{\partial S}\mathbf E\cdot d\mathbf l=-\int_S\frac{\partial\mathbf B}{\partial t}\cdot d\mathbf S$\\[4pt]$\nabla\times\mathbf E=-\partial_t\mathbf B$};
\node[card,text width=6.55cm,draw=purple!65,fill=purple!4] at (4,-1.25) {安培—麦克斯韦定律\\[4pt]$\displaystyle\oint_{\partial S}\mathbf B\cdot d\mathbf l=\mu_0\int_S(\mathbf J+\varepsilon_0\partial_t\mathbf E)\cdot d\mathbf S$\\[4pt]$\nabla\times\mathbf B=\mu_0(\mathbf J+\varepsilon_0\partial_t\mathbf E)$};

\end{tikzpicture}\end{document}
